\section{Introduction}\label{sec:intro}

The inverse Galois problem asks whether every finite group occurs as the
Galois group of a finite extension of $\Q$. The Suzuki group
\[
 N=\Sz(8),\qquad |N|=29120=2^6\cdot5\cdot7\cdot13,
\]
is the smallest of the seven simple groups of order at most $10^8$ on
Zywina's August 2025 list of unresolved cases~\cite{Zywina2025}.
The case of $M_{23}$ on that list has since been resolved by
Huang--Jackson--Lee--Poonen--Pries--Zhang~\cite{HuangEtAl2026}.
We construct a regular realization of $N$, and a specific number-field
realization, by computing an auxiliary cover with group
\[
 G=\Aut(N)=N\rtimes C_3.
\]
Throughout, regular means that the algebraic closure of $\Q$ in the Galois
extension is $\Q$.

\paragraph{The idea.}
For a three-point regular $G$-extension $L/\Q(t)$, the fixed field
$L^N$ is cyclic cubic and ramified at exactly two geometric points.
Its curve has genus zero, and the degree-three map to $\PP^1$ supplies
a rational divisor of odd degree, hence a rational point. Thus
$L^N=\Q(s)$, and $L/\Q(s)$ realizes $N$ regularly
(\cref{cor:threepoint}). This reduces the problem to constructing a
three-point realization of $G$ and identifying its cubic quotient map.
The cover used here has genus $6$ in the degree-$65$ action; its passport
has four Nielsen classes, so rigidity alone does not supply the realization.

The exact data and verification programs form an accompanying proof package,
available at \url{https://github.com/sebastianbiallas/sz8-over-q}.
Let $f(X,t)\in\Z[1/13][X,t]$ be the polynomial supplied there in
\texttt{data/f.json}, monic of degree $65$ in $X$ and of degree $7$ in $t$.
Put
\begin{equation}\label{eq:basechange}
 T(s)=\frac{s^3+6s^2-27s+21}{3(s^3-7s^2+12s-5)}.
\end{equation}

\begin{theorem}\label{thm:main}
The splitting field of $f(X,t)$ over $\Q(t)$ is regular with group
$\Sz(8)\rtimes C_3$. The splitting field of $f(X,T(s))$ over $\Q(s)$ is
regular with group $\Sz(8)$. Moreover,
\[
 \Gal\bigl(f(X,-7/5)/\Q\bigr)\simeq\Sz(8).
\]
Here the Galois group of a polynomial means the Galois group of its
splitting field.
\end{theorem}

A smaller presentation is the monic polynomial $F_3(Y,U)$ of bidegree
$(65,7)$ supplied in \texttt{data/F3\_L3P0P4.json}, with $U=3(t+1)$.
An exact change of generator relates it to $f$ (\cref{sec:smaller}).

\begin{corollary}\label{cor:smaller}
The splitting field of $F_3(Y,3(t+1))$ over $\Q(t)$ is regular with group
$\Sz(8)\rtimes C_3$. Its pullback $F_3(Y,3(T(s)+1))$ has regular splitting
field with group $\Sz(8)$ over $\Q(s)$. The monic integral polynomial
$P_1$ printed in \cref{app:polynomial} defines the degree-$65$ field
obtained by specializing this cover at $s=1$; its Galois group over $\Q$
is $\Sz(8)$ and its signature is $(1,32)$.
\end{corollary}

The auxiliary degree-$65$ cover has genus $6$ and branch-cycle types
\begin{equation}\label{eq:passport}
 \begin{array}{c|ccc}
 \text{branch point}&\zetaT&\zetaT^2&\infty\\ \hline
 \text{cycle type}&1^5 3^{20}&1^5 3^{20}&13^5.
 \end{array}
\end{equation}
The cyclic cubic quotient is branched at the two finite branch points.
Determining its arithmetic form gives \eqref{eq:basechange}.
After this base change the degree-$65$ cover has genus $26$, with three
geometric branch points of type $13^5$.

\begin{figure}[ht]
\centering
\begin{tikzcd}[row sep=2.4em, column sep=2.2em, every label/.append style={font=\footnotesize}]
 & L \arrow[d, dash, "448"]
     \arrow[ddl, dash, bend right=28, "N=\Sz(8)"'] & \\
 & \Q(s,x)\ (26) \arrow[dl, dash, "65"'] \arrow[dr, dash, "C_3"] & \\
 \Q(s)\ (0) \arrow[dr, dash, "C_3"', "t=T(s)"] & & \Q(t,x)\ (6) \arrow[dl, dash, "65"] \\
 & \Q(t)\ (0) &
\end{tikzcd}
\caption{The fields in \cref{thm:main}. Here $L$ is the splitting field of
$f(X,t)$ over $\Q(t)$, with $\Gal(L/\Q(t))=G=\Sz(8)\rtimes C_3$, and $x$ is a
root of $f(X,t)$. Edge labels are degrees or Galois groups; numbers in
parentheses are genera. The fixed field $\Q(s)=L^N$ is rational by
\cref{lem:rationalquotient}, and $t=T(s)$ is \eqref{eq:basechange}.}
\label{fig:tower}
\end{figure}

\paragraph{Contributions and related work.}
The main contribution is the explicit regular realization in
\cref{thm:main}. Two features make the computation practical.
Functions with poles supported over infinity give predictable
coefficient-degree bounds and much smaller equations (\cref{sec:smaller}).
A local discriminant argument removes hundreds of apparent branch points
before monodromy is tracked (\cref{sec:certificate}).

Rational-quotient descent has classical precedents discussed in
\cref{sec:strategy}. Numerical construction of three-point covers follows
the triangle-group approach of Klug--Musty--Schiavone--Voight~\cite{KMSV2014};
see also the survey~\cite{SijslingVoight2014}. Certified path tracking for
Belyi maps has substantial prior work, including Guillemot--Voight's
certification of the Belyi maps in the LMFDB~\cite{GuillemotVoight2026}.
Our emphasis is the arithmetic application and the concrete bounds for
this singular plane model. For the broader arithmetic context of Hurwitz
spaces, see D\`ebes~\cite{Debes2026}.

\paragraph{Proof structure.}
The proof starts from the exact coefficients of $f$ and is independent
of the numerical construction. After the quotient criterion in
\cref{sec:strategy}, \cref{sec:certificate} establishes geometric
monodromy and regularity. \Cref{sec:quotient} identifies the cubic map
from its fibers at infinity and at $t=0$, proving the regular
$\Sz(8)$ realization. \Cref{sec:specialization} completes
\cref{thm:main} with a rational specialization. \Cref{sec:smaller}
then transfers the certificate to $F_3$ by an exact change of generator
and proves \cref{cor:smaller}, including the shorter printed
specialization. \Cref{sec:construction} describes discovery,
and \cref{sec:computation} documents the verification programs and their
exact inputs.
